% Generated by scripts/build-cv.mjs from public portfolio content.
\documentclass[a4paper,10pt]{ltjsarticle}
\usepackage[haranoaji]{luatexja-preset}
% Shared typesetting for the public CVs. Compile with LuaLaTeX.
\usepackage[a4paper,top=17mm,bottom=17mm,left=18mm,right=18mm,footskip=9mm]{geometry}
\usepackage{fontspec}
\setmainfont{LibertinusSerif-Regular.otf}[
  BoldFont=LibertinusSerif-Bold.otf,
  ItalicFont=LibertinusSerif-Italic.otf,
  BoldItalicFont=LibertinusSerif-BoldItalic.otf]
\usepackage{microtype}
\usepackage{titlesec}
\usepackage{enumitem}
\usepackage{needspace}
\usepackage{fancyhdr}
\usepackage{lastpage}
\usepackage{hyperref}
\hypersetup{hidelinks,unicode=true,pdfauthor={Taichi Uchida},pdfsubject={Public curriculum vitae}}
\setlength{\parindent}{0pt}
\setlength{\parskip}{0pt}
\setlength{\emergencystretch}{2em}
\linespread{1.05}
\raggedbottom
\pagestyle{fancy}
\fancyhf{}
\fancyfoot[L]{\footnotesize Taichi Uchida\quad\textbar\quad Curriculum Vitae}
\fancyfoot[R]{\footnotesize\thepage\,/\,\pageref*{LastPage}}
\renewcommand{\headrulewidth}{0pt}
\renewcommand{\footrulewidth}{0.3pt}
\titleformat{\section}{\large\bfseries}{}{0pt}{}[\vspace{1pt}\titlerule]
\titlespacing*{\section}{0pt}{10pt}{6pt}
\setlist[itemize]{leftmargin=1.2em,itemsep=1pt,topsep=2pt,parsep=0pt}
\newcommand{\cvsection}[1]{\Needspace{5\baselineskip}\section*{#1}}
% A fixed date column and an unbroken entry keep periods and descriptions aligned.
\newcommand{\cventry}[3]{%
  \par\addvspace{5pt}\noindent
  \begin{minipage}[t]{28mm}\small\raggedright #1\end{minipage}\hfill
  \begin{minipage}[t]{\dimexpr\textwidth-33mm\relax}
    \textbf{#2}\par\vspace{1.5pt}#3
  \end{minipage}\par
}
\newcommand{\cvcompact}[3]{%
  \par\addvspace{4pt}\noindent
  \begin{minipage}[t]{28mm}\small\raggedright #1\end{minipage}\hfill
  \begin{minipage}[t]{\dimexpr\textwidth-33mm\relax}
    \textbf{#2}\if\relax\detokenize{#3}\relax\else\enspace #3\fi
  \end{minipage}\par
}

\hypersetup{pdftitle={内田 大智 - Curriculum Vitae}}
\begin{document}
{\fontsize{25}{30}\selectfont\bfseries 内田 大智}\hspace{5mm}{\large Taichi Uchida}\hfill{\small 2026年9月更新}\par\vspace{5pt}
京都大学 工学部物理工学科 宇宙基礎工学コースの学部4年生。ミッション設計、軌道最適化、ロボティクスを研究しています。\par\vspace{4pt}
{\small\href{mailto:uchida.taichi.84f@st.kyoto-u.ac.jp}{uchida.taichi.84f@st.kyoto-u.ac.jp}\quad\textbar\quad\href{https://taichi-u.github.io/}{taichi-u.github.io}}\par\vspace{4pt}
\cvsection{学歴}
\cventry{2027.04}{総合研究大学院大学・進学予定}{先端学術院 先端学術専攻 宇宙科学コース・5年一貫博士課程。尾崎直哉研究室（Mission Design Lab）。}
\cventry{2023.04 - 2027.03}{京都大学}{工学部物理工学科 宇宙基礎工学コース・学部4年生。2027年3月卒業見込。}
\cventry{2020 - 2023}{聖光学院高等学校}{神奈川県横浜市}
\cventry{2018 - 2020}{Pathways School Gurgaon}{IB Middle Years Programme・インド ハリヤナ州}
\cventry{2017 - 2018}{聖光学院中学校}{神奈川県横浜市}

\cvsection{研究・開発}
\cventry{2025.08 -}{外惑星探査のミッション設計と軌道最適化}{JAXA / ISAS / テクニカル研修生・尾崎直哉研究室. 超小型外惑星探査機による多点同時観測のミッション設計と、Julia・DDPを用いた低推力軌道最適化。}
\cventry{2026.04 -}{Transformerを用いたモデル予測制御}{KYOTO UNIVERSITY / 学部研究・藤本健治研究室（制御工学分野）. Transformerで同定したモデルと非線形モデル予測制御を組み合わせ、ラジコンカーの実時間自動追従を実機に実装。}
\cventry{2025.10 -}{月面ロボットの強化学習と制御}{KYOTO UNIVERSITY / 研究補助・学習機械制御論分野（森本研究室）. 月面探査用の多脚モジュールロボットを対象に、強化学習・逆運動学・Sim-to-Real環境を研究。}
\cventry{2025.02 - 2025.04}{7自由度ロボットアームの制御}{DFKI · BREMEN / 客員研究員・ロボティクス・イノベーション・センター. 宇宙デブリ捕獲を背景に、7自由度アームの物理シミュレーション環境とインピーダンス・位置制御器を開発。}
\cventry{2025 -}{陸上競技の記録管理Webアプリ}{KYOTO UNIVERSITY · TRACK \& FIELD / 開発・システム管理. 京都大学医学部陸上部に向けて、React・Firebaseによる競技記録管理アプリを設計・開発・運用。}
\cventry{2025.03 -}{ロシア語の文法誤り訂正AI}{TLEEZ INC. / 自然言語処理エンジニア インターン. ロシア語学習を支援する文法誤り訂正AIと、学習データ生成に取り組むエンジニアインターン。}
\cventry{2024.08 -}{微生物群集のメタゲノム解析}{GEORGIA TECH / 客員研究員・DiChristina研究室. 重金属汚染物質の安定化に関わる微生物群集を、メタゲノムデータと統計解析から調べる。}
\cventry{2020.12 - 2023.03}{岩塩中の好塩性古細菌の単離・同定}{TOKYO TECH · ELSI / 研究インターン・藤島皓介研究室. 世界各地の岩塩から好塩性古細菌を単離し、DNA解析と分子系統解析を実施。}
\newpage
\cvsection{その他の研究}
\cventry{2023.09 - 2024.08}{微生物代謝研究・京都大学}{跡見研究室で、水素・酸素の利用に関する微生物代謝のウェットラボ研究。}
\cventry{2021.08 - 2022.02}{隕石分析・聖光学院}{電子顕微鏡（SEM）による組成分析。英国の高校・大学教員との国際共同研究。}
\cventry{2020.08 - 2021.08}{人工重力機構・聖光学院}{新規機構を試作し、Arduinoを用いた力センサーを設計・実装して評価。}
\cventry{2021}{知能化水ロケット}{コンピュータ制御機能と自動パラシュート展開機構を備えたロケットを開発。}

\cvsection{発表・執筆}
\cventry{2026.08}{GTOからGEOへの低推力軌道遷移におけるHDDPを用いた軌道最適化}{70th UKAREN. 論文投稿.}
\cventry{2026.03}{\href{https://pparc.tohoku.ac.jp/sympo/sps/archive/sps2026/}{海外大型ミッション相乗りを想定した超小型外惑星探査機の多点同時観測のミッション設計}}{SPS 2026. 口頭発表.}
\cventry{2023.11}{日本における宇宙医学教育の現状と展望}{69th JSAEM. 口頭・ポスター発表.}
\cventry{2022.06}{世界各地の岩塩からの極限環境微生物の分離同定}{NASA × AGI × High School Students. オンライン発表.}
\cventry{2022.05}{世界各地の岩塩からの高度好塩菌の単離と同定}{JpGU 2022. ポスター発表・優秀賞.}
\cventry{2022.03}{極限環境微生物研究}{Japanese Society of Scientific Fisheries. オンライン発表.}
\cventry{2022.01}{電子顕微鏡による隕石組成分析}{International Research for School. オンライン発表.}

\cvsection{受賞・採択}
\cvcompact{2025}{Startup Weekend Space Osaka・準優勝}{チームリーダーとして54時間で事業案とプロトタイプを開発。}
\cvcompact{2024-25}{中谷財団 海外留学奨学金}{ジョージア工科大学・DFKIへの研究派遣（全額給付）。}
\cvcompact{2022}{日本地球惑星科学連合 高校生セッション優秀賞}{好塩性古細菌の研究。}
\cvcompact{2022}{日本水産学会 高校生優秀発表賞}{極限環境微生物の研究。}
\cvcompact{2022}{SDGs探究アワード・日本旅行賞}{Seiko SDGs Food Ambassadorの創設・運営。}
\cvcompact{2021}{知能化水ロケット競技会 ミッション賞}{競技会の全ミッションを達成。}
\cvcompact{2021}{SSH生徒研究発表会 優秀賞}{人工重力機構の研究。}
\cvcompact{2021}{ゆりかもめリレーマラソン・総合優勝}{チームでの総合優勝。}
\cvcompact{2019}{ISSO全国大会・3種目優勝}{インドで800m・1500m・3000mの3種目優勝。}
\newpage
\cvsection{活動歴}
\cventry{2026.09.15}{OpenAI 大学生向け 夏の Codex 開発祭 2026}{日陰を通る最短経路を探索する地図サービス「WalkIsFun」を3人チームで開発。PMを担当しました。}
\cventry{2026.03}{JENESYS2025・日韓青少年交流訪韓団}{外務省・日韓文化交流基金のプログラムに派遣団員として参加。}
\cventry{2025.05}{京都・国際音楽学生フェスティバルでの通訳}{京都・国際音楽学生フェスティバルで、海外音楽家の英日通訳とアテンドを担当。}
\cventry{2023 -}{鉄緑会 数学講師}{鉄緑会で難関大学志望の高校生に数学を指導。}
\cventry{2021}{日本学生宇宙フォーラム}{100名以上が参加するフォーラムを創設・企画運営。}
\cventry{2021}{聖光SDGs食アンバサダー}{文化祭での地域産品の販売を通じて生産者・観光を支援する活動を創設。}
\cventry{中学時代}{「横浜バナナ」プロジェクト}{プロジェクトを創設・運営し、2日間で3,800本のチョコバナナを販売。}
\cventry{2025 / 2023}{宇宙医学・Moon Village}{第71回日本宇宙航空環境医学会・第7回Global Moon Village Workshopの運営、大学生向け宇宙医学ツアーの企画。}
\cventry{2012 -}{稲作と少林寺拳法}{有機稲作への継続的な参加。少林寺拳法では全国中学生大会に出場。}
\cventry{2021 / 2019}{科学と国際交流}{科学部サミットを創設。日印学生交流の企画や、インドでの書道体験イベントを運営。}
\cventry{2021}{議員インターン}{ICASプログラムで東京都議会の議員インターンに参加。}

\cvsection{スキル・資格}
\cvcompact{}{プログラミング}{Python · Julia · C++ · MATLAB · R · Go}
\cvcompact{}{制御・ロボティクス}{DDP · NMPC · Drake · Pinocchio · MuJoCo · Isaac Lab · CasADi · ROS2}
\cvcompact{}{機械学習・Web開発}{PyTorch · React · Next.js · Firebase · Docker · Linux · Git}
\cvcompact{}{生命科学}{PCR · DNA extraction · PhyML · Metagenomics}
\cvcompact{}{言語}{日本語（母語）・英語（TOEFL iBT 100／英検1級）・ロシア語（学習中）}
\cvcompact{}{資格}{モデルロケットライセンス4級・NAUIオープンウォーターダイバー・普通自動車免許（AT限定）}
\end{document}
